Der Begriff \emph{Faschismus} wird im öffentlichen Diskurs fast ausschließlich mit dem Nationalsozialismus, Hitler und dem Holocaust gleichgesetzt. Diese Gleichsetzung ist historisch falsch.

Der italienische Faschismus unter Benito Mussolini war:
\begin{itemize}
    \item \textbf{Nicht von Judenhass getrieben}: Mussolini bezeichnete sich selbst als nicht judenfeindlich. Erst 1938, unter Hitlers Druck, wurden Rassengesetze erlassen.
    \item \textbf{Nationalistisch}: Der Faschismus sah den Staat als alles und das Individuum als nichts – \emph{Tutto nello Stato, niente fuori dallo Stato, nulla contro lo Stato}.
    \item \textbf{Antidemokratisch}: Er zerschlug die Demokratie, verbot Gewerkschaften und errichtete eine Diktatur.
\end{itemize}

Die entscheidende historische Einsicht ist jedoch eine andere: Der Faschismus entstand nicht aus dem Nichts, sondern aus der \emph{Angst des Kapitals vor der Demokratie}. Die Industriellen und Großgrundbesitzer Italiens fürchteten die sozialistische Bewegung, die nach dem Ersten Weltkrieg in Italien stark war (Biennio Rosso 1919–1920). Sie finanzierten Mussolini, um die Arbeiterbewegung zu zerschlagen – und wurden später selbst von ihm gefressen.

In Deutschland wiederholte sich dieses Muster: Die Industriellen (Thyssen, Krupp, IG Farben, Deutsche Bank) finanzierten die NSDAP, um die KPD und SPD zu zerschlagen. Sie glaubten, Hitler kontrollieren zu können – und wurden von ihm gefressen.

Die strukturelle Definition des Faschismus lautet daher:

\begin{quote}
\emph{Faschismus ist die Steuerung der Politik durch das Kapital, und Faschismus ist immer antidemokratisch.}
\end{quote}

Diese Definition löst den Begriff von der historischen Erscheinungsform (Hakenkreuz, SA, Diktator) und macht ihn für die Gegenwart anwendbar – unabhängig davon, ob die faschistische Struktur mit Gewalt oder mit medialer Steuerung, Lobbyismus und Parteispenden arbeitet.